% 问题二三线表。需要 \usepackage{booktabs}、\usepackage{siunitx}。
% 数值来源：results/group_statistics.csv、paired_layout_comparisons.csv、paired_dop_comparison.csv。

\begin{table}[htbp]
  \centering
  \caption{不同观测布局的定位质量比较（30次配对重复）}
  \label{tab:p2-layout-comparison}
  \small
  \resizebox{\textwidth}{!}{%
  \begin{tabular}{lrrrrrr}
    \toprule
    布局方法 & 面积/\si{\metre\squared} & 直径/\si{\metre} & 覆盖半径/\si{\metre}
      & 定位误差/\si{\metre} & $\widetilde\gamma_{\min}$/\si{\degree} & DOP \\
    \midrule
    随机布局       & 638.527 & 47.470 & 23.941 & 11.389 & 13.0 & 1.231 \\
    均匀圆周布局   & 520.174 & 32.457 & \textbf{16.318} & 8.271 & 0.0 & 1.000 \\
    本文$\Omega$贪心布局 & \textbf{462.886} & \textbf{31.856} & 16.453 & \textbf{8.182} & 45.0 & 1.000 \\
    DOP诊断基线    & 574.535 & 34.762 & 18.111 & 10.232 & 45.0 & 1.000 \\
    \bottomrule
  \end{tabular}}

  \vspace{1mm}
  \parbox{0.96\linewidth}{\footnotesize 注：表内为30次重复的均值；面积为问题一网格内外包中点，直径为内包近似，覆盖半径为安全外包圆半径。DOP仅作诊断，不是最终优化目标。均匀圆周布局含反向共线点对，故最小条件交会角为0°，但其他独立方向仍可使整体几何良好。}
\end{table}

\begin{table}[htbp]
  \centering
  \caption{各布局相对随机布局的配对差}
  \label{tab:p2-paired-layout}
  \small
  \resizebox{\textwidth}{!}{%
  \begin{tabular}{lrrrr}
    \toprule
    布局方法 & 面积差/\si{\metre\squared} & 面积差95\%自举区间
      & 直径差/\si{\metre} & 定位误差差/\si{\metre} \\
    \midrule
    均匀圆周布局 & -118.353 & $[-245.005,\ 6.289]$ & -15.012 & -3.118 \\
    本文$\Omega$贪心布局 & -175.641 & $[-329.210,\ -23.165]$ & -15.613 & -3.207 \\
    DOP诊断基线 & -63.992 & $[-232.959,\ 94.206]$ & -12.708 & -1.157 \\
    \bottomrule
  \end{tabular}}

  \vspace{1mm}
  \parbox{0.96\linewidth}{\footnotesize 注：差值定义为“相应方法减随机布局”，负值表示指标减小。自举区间针对重复场景均值，不是干扰源位置的概率置信域。}
\end{table}

\begin{table}[htbp]
  \centering
  \caption{本文布局下检测点数量对定位质量的影响}
  \label{tab:p2-number}
  \small
  \begin{tabular}{rrrrr}
    \toprule
    检测点数$n$ & 面积/\si{\metre\squared} & 直径/\si{\metre}
      & 覆盖半径/\si{\metre} & 定位误差/\si{\metre} \\
    \midrule
    2 & 1230.922 & 49.744 & 24.969 & 14.019 \\
    3 & 746.004  & 45.130 & 22.671 & 10.921 \\
    4 & 462.886  & 31.856 & 16.453 & 8.182 \\
    5 & 336.061  & 29.684 & 15.146 & 7.577 \\
    6 & 237.768  & 25.468 & 12.963 & 5.905 \\
    \bottomrule
  \end{tabular}
\end{table}

\begin{table}[htbp]
  \centering
  \caption{两站原始交会角对区域定位结果的影响}
  \label{tab:p2-angle}
  \small
  \begin{tabular}{rrrrr}
    \toprule
    原始交会角$\alpha$/\si{\degree} & 面积/\si{\metre\squared}
      & 直径/\si{\metre} & 覆盖半径/\si{\metre} & 定位误差/\si{\metre} \\
    \midrule
    30  & 2536.352 & 138.126 & 69.255 & 27.394 \\
    45  & 1766.752 & 92.453  & 46.363 & 19.452 \\
    60  & 1432.368 & 70.418  & 35.322 & 16.056 \\
    90  & \textbf{1230.922} & \textbf{49.744} & \textbf{24.969} & \textbf{14.019} \\
    120 & 1413.802 & 69.516  & 34.881 & 15.799 \\
    \bottomrule
  \end{tabular}

  \vspace{1mm}
  \parbox{0.96\linewidth}{\footnotesize 注：此表横轴是两站原始夹角$\alpha\in[0,180^\circ]$；多站布局诊断采用$\widetilde\gamma=\min(\gamma,180^\circ-\gamma)$，故原始120°对应条件角60°。}
\end{table}

\begin{table}[htbp]
  \centering
  \caption{本文$\Omega$贪心布局与DOP基线的直接配对比较}
  \label{tab:p2-dop-direct}
  \small
  \begin{tabular}{lrr}
    \toprule
    指标（本文减DOP） & 配对均值差 & 95\%自举区间 \\
    \midrule
    面积/\si{\metre\squared} & -111.649 & $[-232.065,\ -0.351]$ \\
    直径/\si{\metre} & -2.905 & $[-5.792,\ -0.354]$ \\
    覆盖半径/\si{\metre} & -1.658 & $[-3.268,\ -0.227]$ \\
    定位误差/\si{\metre} & -2.050 & $[-3.995,\ -0.275]$ \\
    \bottomrule
  \end{tabular}
\end{table}
